
\documentclass[../../../physics-standard_answers_all.tex]{subfiles}

\begin{document}

\setlength{\abovedisplayskip}{2mm}
\setlength{\belowdisplayskip}{1mm}

\begin{enumerate}

    \item　
        \par {\centering
            \includegraphics
                {\subfix{fig/fig_em_ef_05/fig_em_ef_05_01.pdf}}
        \par}
        \vspace{0.2\baselineskip}

    \item $x$軸上での電位は$\phi = ax^2$である．
    $x>0$の領域での電場は \:\,\uwave{$x$軸の負の向き} であり，強さは，
        %
        \begin{align*}
            \left| E_x \right| = \left| \lim_{\Delta x \to 0}
            \frac{\Delta \phi}{\Delta x} \right| = 2ax \;.
        \end{align*}
        %
    $x=1 \, \mathrm{m}$の位置では，
        \begin{align*}
            \left| E_x \right| = \uwave{20 \, \mathrm{V/m}} \;.
        \end{align*}

    \item $y$軸上での電位は$\phi = -ay^2$である．
    位置エネルギーが最大（運動エネルギーが最小）となる原点Oを通過できればよい．
    原点に到達したときの速さを$u$とすれば，力学的エネルギー保存則より，
        %
        \begin{align*}
            \frac{1}{2}mu^2 + q \cdot 0 = \frac{1}{2}m{v_0}^2 + q(-ah^2) \;.
        \end{align*}
        %
        よって，$ \displaystyle \frac{1}{2}mu^2 = \frac{1}{2}m{v_0}^2 -qah^2 > 0$となるために，
        %
        \begin{align*}
            \uwave{v_0 > h \sqrt{\frac{2qa}{m}}} \;.
        \end{align*}

\end{enumerate}

\end{document}
